\documentclass[11pt]{amsart}
\usepackage{amsmath,amssymb,amsthm}
\usepackage{graphicx}
\usepackage[hidelinks]{hyperref}

\newtheorem{theorem}{Theorem}
\newtheorem{lemma}{Lemma}
\newcommand{\M}{\mathcal M}
\newcommand{\Tr}{\operatorname{Tr}}
\newcommand{\supp}{\operatorname{supp}}

\title[Stability for a uniformly atomic type-I subcase]
{A uniformly atomic type-I subcase of a stability problem for noncommutative symmetric spaces}
\author{fa\_banach\_001, agent\_lane\_07}
\date{2026-07-19}

\begin{document}
\maketitle

\section*{Status and source}

This packet gives a candidate partial positive answer to Problem 9 in
Czerwinska--Kaminska \cite{CK2017}, arXiv:1704.02033, PDF page 55.  The source
asks:
\begin{quote}
``Assume $E$ is stable and $\M$ is of type I. Is the space $E(\M,\tau)$
stable?''
\end{quote}
The result below answers this affirmatively for countable atomic type-I
algebras whose positive factor trace weights are uniformly bounded below.  The
weights need not be commensurable.  It does not settle the diffuse-center case
or atomic weights tending to zero.

The original page also records the Arazy--Raynaud theorem used below and is
shown here at full readable width.

\begin{center}
\includegraphics[width=0.98\textwidth]{figures/open_problem_crop.png}
\end{center}

\section*{Definitions}

A metric space $(X,d)$ is \emph{stable} if for every two bounded sequences
$(x_m)$ and $(y_n)$ and every two free ultrafilters $\mathcal U,\mathcal V$ on
$\mathbb N$,
\[
 \lim_{m,\mathcal U}\lim_{n,\mathcal V}d(x_m,y_n)
 =
 \lim_{n,\mathcal V}\lim_{m,\mathcal U}d(x_m,y_n).
\]
A Banach space is stable when its norm metric is stable.  Stability is
inherited by isometric subspaces.

For a symmetric Banach sequence space $F$, its unitary matrix ideal is
\[
 C_F=\{T\in\mathcal K(K):s(T)=(s_j(T))_{j\geq1}\in F\},
 \qquad \|T\|_{C_F}=\|s(T)\|_F,
\]
where $s_j(T)$ are the singular values in decreasing order.  The
Arazy--Raynaud theorem says that a separable symmetric sequence space $F$ is
stable if and only if $C_F$ is stable \cite{Arazy1983,CK2017}.

For a semifinite von Neumann algebra $(\M,\tau)$ and a symmetric Banach
function space $E$, the noncommutative symmetric space $E(\M,\tau)$ consists
of the $\tau$-measurable operators $x$ for which the generalized singular-value
function $\mu_\tau(x)$ belongs to $E$, with
$\|x\|_{E(\M,\tau)}=\|\mu_\tau(x)\|_E$.

\section*{Proof Intuition}

The obstruction in the full problem is that a general type-I trace need not
look like ordinary counting measure on one Hilbert space.  In the atomic
commensurable case it does, after an exact amplification.  If the trace on the
$n$th factor is $\delta k_n\Tr$, tensor that factor with the identity on
$\mathbb C^{k_n}$.  Every singular value is then repeated $k_n$ times, so the
ordinary spectral counting function of the amplified operator is exactly
$1/\delta$ times its original $\tau$-distribution function.

Equal intervals of length $\delta$ cut a symmetric sequence space $F_\delta$
out of $E$.  The amplification identifies $E(\M,\tau)$ isometrically with a
closed subspace of $C_{F_\delta}$.  Because $F_\delta$ is a closed subspace of
the stable space $E$, it is stable; Arazy--Raynaud then makes
$C_{F_\delta}$ stable.  If the original weights merely have a positive lower
bound, round them upward to multiples of a common, very small $\delta$.  This
changes every norm by a factor at most $1+\varepsilon$.  Stability survives the
limit of these arbitrarily small distortions.  Diffuse centers and weights
tending to zero remain open after this argument.

\newpage
\section*{Partial theorem}

\begin{theorem}[Commensurable atomic case]\label{thm:commensurable}
Let $E$ be a stable, separable symmetric Banach function space on
$[0,\infty)$ with order-continuous norm.  Let
\[
 \M=\mathop{\bigoplus}^{\ell_\infty}_{n\geq1}B(H_n),
\]
where every $H_n$ is a separable Hilbert space.  Suppose that, for some
$\delta>0$ and positive integers $k_n$, the faithful normal semifinite trace
is
\[
 \tau((x_n))=\delta\sum_{n\geq1}k_n\Tr_{H_n}(x_n),
 \qquad (x_n)\in\M_+.
\]
Then $E(\M,\tau)$ is stable.
\end{theorem}

\begin{proof}
Put $I_j=[(j-1)\delta,j\delta)$ and let $X_\delta$ be the closed linear span
in $E$ of the pairwise disjoint indicators $\mathbf1_{I_j}$.  Define the
coefficient sequence space
\[
 F_\delta=\left\{a=(a_j):
     \sum_{j\geq1}a_j\mathbf1_{I_j}\text{ converges in }E\right\},
 \qquad
 \|a\|_{F_\delta}
 =\left\|\sum_{j\geq1}a_j\mathbf1_{I_j}\right\|_E.
\]
Because the intervals are equimeasurable and $E$ is symmetric,
$F_\delta$ is a symmetric Banach sequence space.  Its step-function map is a
linear isometry onto $X_\delta$.  Thus $F_\delta$ is stable, since
$X_\delta$ is a closed subspace of the stable Banach space $E$; it is also
separable.  The Arazy--Raynaud theorem therefore implies that
$C_{F_\delta}$ is stable.

Let
\[
 K=\bigoplus_{n\geq1}(H_n\otimes\mathbb C^{k_n})
 \quad\text{and}\quad
 J((x_n))=\bigoplus_{n\geq1}(x_n\otimes I_{k_n}).
\]
Functional calculus commutes with this amplification.  Consequently, for
every $x$ affiliated with $\M$ and every $\lambda\geq0$,
\begin{align*}
 \Tr_K\!\left(\mathbf1_{(\lambda,\infty)}(|Jx|)\right)
 &=\sum_{n\geq1}k_n
   \Tr_{H_n}\!\left(\mathbf1_{(\lambda,\infty)}(|x_n|)\right)\\
 &=\delta^{-1}\tau\!\left(
   \mathbf1_{(\lambda,\infty)}(|x|)\right).
\end{align*}
The equality is valid in $[0,\infty]$.  Taking generalized inverses of these
distribution functions gives
\[
 \mu_\tau(x)(t)=\mu_{\Tr_K}(Jx)(t/\delta),\qquad t\geq0.
\]
Set
\[
 a_j=\mu_{\Tr_K}(Jx)(j-1),\qquad j\geq1.
\]
These are the decreasing spectral values of the amplification (with extended
values allowed at this point).  The scaling identity says exactly
\begin{equation}\label{eq:step}
 \mu_\tau(x)
 =\sum_{j\geq1}a_j\mathbf1_{I_j}.
\end{equation}

Now take $x\in E(\M,\tau)$.  Order continuity of $E$ makes the finite
truncations of the right side of \eqref{eq:step} converge to
$\mu_\tau(x)$ in $E$.  Hence $a=(a_j)\in F_\delta$.  Since a separable
symmetric sequence space is contained in $c_0$, the sequence tends to zero.
Thus $Jx$ is a compact bounded operator, $a=s(Jx)$ is its usual singular-value
sequence, and $Jx\in C_{F_\delta}$.  Moreover,
\[
 \|x\|_{E(\M,\tau)}
 =\|\mu_\tau(x)\|_E
 =\|s(Jx)\|_{F_\delta}
 =\|Jx\|_{C_{F_\delta}}.
\]
Thus $J$ is a linear isometry from $E(\M,\tau)$ into
$C_{F_\delta}$.  Its range is closed because the domain is complete.
Stability of $C_{F_\delta}$ passes to this closed subspace, proving the
claim.
\end{proof}

\begin{lemma}[Almost-isometric stability test]\label{lem:almost}
Let $X$ be a Banach space.  Suppose that for every $c>1$ there are a stable
Banach space $Y_c$ and a linear map $T_c:X\to Y_c$ such that
\[
 \|z\|_X\leq\|T_cz\|_{Y_c}\leq c\|z\|_X
 \qquad(z\in X).
\]
Then $X$ is stable.
\end{lemma}

\begin{proof}
For bounded sequences $(x_n),(y_m)$ and ultrafilters $\mathcal U,\mathcal V$,
let $A$ and $B$ be the two iterated limits of $\|x_n+y_m\|_X$ in opposite
orders.  Stability of $Y_c$ says that the corresponding two iterated limits of
$\|T_cx_n+T_cy_m\|_{Y_c}$ have one common value $L_c$.  The distortion bounds
give
\[
 A\leq L_c\leq cA,
 \qquad B\leq L_c\leq cB.
\]
Hence $A\leq cB$ and $B\leq cA$ for every $c>1$.  Letting $c\downarrow1$
gives $A=B$.
\end{proof}

\begin{theorem}[Uniformly atomic case]\label{thm:uniform}
Let $E$ be as in Theorem~\ref{thm:commensurable}, and let
\[
 \M=\mathop{\bigoplus}^{\ell_\infty}_{n\geq1}B(H_n),
 \qquad
 \tau((x_n))=\sum_{n\geq1}w_n\Tr_{H_n}(x_n),
\]
where $w_n>0$ and $w_*:=\inf_n w_n>0$.  Then $E(\M,\tau)$ is stable.  No
commensurability assumption on $(w_n)$ is required.
\end{theorem}

\begin{proof}
Fix $c>1$.  Choose $\delta>0$ with $\delta<(c-1)w_*$, put
$k_n=\lceil w_n/\delta\rceil$, and set $v_n=\delta k_n$.  Then
\[
 w_n\leq v_n<cw_n\qquad(n\geq1).
\]
Let $\tau_v=\sum_n v_n\Tr_{H_n}$.  For every affiliated $x$ and every
$\lambda\geq0$, the two distribution functions satisfy
\[
 d_{\tau}(\lambda;x)\leq d_{\tau_v}(\lambda;x)
 \leq c\,d_{\tau}(\lambda;x).
\]
Taking generalized inverses yields
\[
 \mu_\tau(x)(t)\leq\mu_{\tau_v}(x)(t)
 \leq \mu_\tau(x)(t/c)=:D_c\mu_\tau(x)(t).
\]
The stable symmetric space $E$ is order continuous and hence fully symmetric
in the source framework.  For decreasing $f\geq0$,
\[
 \int_0^tD_cf(s)\,ds
 =c\int_0^{t/c}f(u)\,du
 \leq c\int_0^t f(u)\,du,
\]
so $D_cf\prec\!\prec cf$ and $\|D_cf\|_E\leq c\|f\|_E$.  Consequently
\[
 \|x\|_{E(\M,\tau)}
 \leq\|x\|_{E(\M,\tau_v)}
 \leq c\|x\|_{E(\M,\tau)}.
\]
The identity is therefore a distortion-$c$ linear embedding of
$E(\M,\tau)$ into $E(\M,\tau_v)$.  The latter is stable by
Theorem~\ref{thm:commensurable}.  Lemma~\ref{lem:almost}, applied for every
$c>1$, proves the claim.
\end{proof}

\section*{Verifier report}

The proof was checked independently along the following chain.
\begin{enumerate}
\item The equal-step coefficient space is not treated as an arbitrary set of
step functions: it is defined through the closed span of the equal interval
indicators.  This makes completeness and the isometric embedding into $E$
explicit.
\item At a spectral threshold $\lambda$, tensoring the $n$th block with
$I_{k_n}$ multiplies its spectral rank by exactly $k_n$.  This proves the
distribution identity before any norm argument and avoids an implicit
reordering assumption.
\item Generalized inversion scales the distribution variable by $\delta$,
which yields intervals of length $\delta$ in \eqref{eq:step}; no missing
factor remains in the norm equality.
\item Order continuity is used exactly once, to pass from finite singular-value
steps to the full function in $E$.  Separability of $F_\delta$ then guarantees
compactness of the amplification.
\item A reusable finite check in
\path{code/check_finite_distribution.py} tested six block families at
nine thresholds and returned \texttt{PASS}.  It checks only the elementary
integer-multiplicity identity and is not part of the proof.
\item For the upgraded theorem, the inequalities between the two distribution
functions were inverted explicitly, and the only dilation estimate used is
$D_cf\prec\!\prec cf$, valid because stable $E$ is order continuous and hence
fully symmetric in the source framework.
\end{enumerate}

No unproved mathematical dependency remains inside the stated theorem.  The
external dependency is the cited Arazy--Raynaud stability theorem for unitary
matrix ideals.  Recommended human-review focus: confirm the source convention
under which stability of a symmetric function space supplies the stated
separability/order-continuity setting, and recheck the generalized-inverse
scaling in \eqref{eq:step} and in the trace-perturbation argument of
Theorem~\ref{thm:uniform}.

\section*{Limitations and novelty status}

The upgraded theorem does not address type-I algebras with diffuse center or a
countable atomic direct sum whose positive trace weights tend to zero.  In
those cases no single commensurable trace gives a uniformly small relative
perturbation.  Local or finite-dimensional approximation is insufficient by
itself because stability is not merely a finite-representability property.

Cheap run indexes were searched for arXiv:1704.02033, the exact title, and the
core stability/type-I/unitary-ideal terms.  A bounded web search on 2026-07-19
using the exact question, atomic/commensurable variants, and uniformly positive
trace weights found the source and the classical sequence-ideal theorem but no
later resolution or duplicate of this trace-perturbation reduction.  This is a
bounded novelty check, not an exhaustive claim of priority.  The status is therefore
\texttt{candidate\_partial\_likely\_valid}, pending human review.

\begin{thebibliography}{3}

\bibitem{CK2017}
M. M. Czerwinska and A. Kaminska,
\emph{Geometric properties of noncommutative symmetric spaces of measurable
operators and unitary matrix ideals}, arXiv:1704.02033 (2017), Problem 9,
PDF p.~55.

\bibitem{Arazy1983}
J. Arazy,
\emph{On stability of unitary matrix spaces},
Proc. Amer. Math. Soc. 87 (1983), 317--321.

\bibitem{Raynaud1983}
Y. Raynaud,
\emph{Stabilite des espaces $C_E$},
Seminaire de Geometrie des Espaces de Banach, Universite Paris 7, 1982--1983.

\end{thebibliography}

\end{document}
